\subsection{自同态的谱}

定义 1.　设 E 为拓扑向量空间，u 为 E 的自同态。我们把 u 相对于幺代数 \(\mathcal{L}(\mathrm{E})\) 的谱称为 u 的谱，记作 Sp(u)。

设 E 为拓扑向量空间，\(u \in \mathcal{L}(\mathrm{E})\) 。u 的谱是由满足以下条件的复数 \(\lambda\) 组成的集合：\(u - \lambda1_{E}\) 不是 E 的自同构。若 E 可度量化且完备，则它也是由满足以下条件的复数 \(\lambda\) 组成的集合：\(u - \lambda1_{E}\) 不是双射（EVT，I，第 19 页，推论 1）。

u 的每个特征值都属于 u 的谱，但反之一般不成立。

在本节余下部分，我们将限于研究 E 为 Banach 空间情形下的谱概念。

引理 1.　设 E 为 Banach 空间，u 为 E 的自同态。设 \((\mathrm{E}_{i})_{i\in\mathrm{I}}\) 为 E 的在 u 下稳定的闭子空间的有限族，满足 \(E = \bigoplus_{i\in I} E_{i}\) 。对每个 \(i \in I\) ，用 \(u_{i}\) 表示 u 在 \(E_{i}\) 上诱导的自同态。我们有 \(\mathrm{Sp}(u) = \bigcup_{i\in I} \mathrm{Sp}(u_{i})\) ，并且对每个 \(f \in \mathcal{O}(\mathrm{Sp}(u))\) ，自同态 \(f(u)\) 使诸空间 \(E_{i}\) 保持稳定，且 \(f(u)\) 与 \(f(u_{i})\) 在 \(E_{i}\) 上一致。

自同态 \(u\) 是同构当且仅当 \(u_i\) 对每个 \(i \in \mathrm{I}\) 都是同构。把这性质应用于 \(u - \lambda 1_{\mathrm{E}}\) ，便推出 \(\operatorname{Sp}(u)\) 是诸集合 \(\operatorname{Sp}(u_i)\) （\(i \in \mathrm{I}\)）的并。

设 \(f \in \mathcal{O}(\mathrm{Sp}(u))\) 。E 的自同态 \(f(u)\) 属于 \(u\) 在 \(\mathcal{L}(\mathrm{E})\) 中的双交换子（I，第 74 页，定理 5），因此它与诸投影 \(p_i\) 交换，从而它使诸空间 \(\mathrm{E}_i\) 保持稳定。考察连续保单位元态射 \(\varpi\) ，它由乘积代数 \(\prod_{i \in \mathrm{I}} \mathcal{L}(\mathrm{E}_i)\) 到 \(\mathcal{L}(\mathrm{E})\) 中的对应 \((v_i)_{i \in \mathrm{I}} \mapsto \bigoplus_i v_i\) 所定义。它把族 \((u_i)\) 映为 \(u\) 。于是 I，第 75 页，命题 7 蕴涵 \(\mathrm{Sp}(u) \subset \bigcup_{i \in \mathrm{I}} \mathrm{Sp}(u_i)\) 以及

\[
f (u) = f (\varpi ((u _ {i}) _ {i \in \mathrm{I}})) = \varpi ((f (u _ {i})) _ {i \in \mathrm{I}}) = \bigoplus_ {i \in \mathrm{I}} f (u _ {i}),
\]

至此引理证毕。

命题 1.　设 E 为复 Banach 空间，u 为 E 的自同态。设 \(\lambda \in \mathbf{C}\) ，\(f \in \mathcal{O}(\mathrm{Sp}(u))\) 。我们有

\[
\operatorname{Ker} (u - \lambda 1 _ {\mathrm{E}}) \subset \operatorname{Ker} (f (u) - f (\lambda) 1 _ {\mathrm{E}}).
\]

设 \(x \in E\) 非零。设 A 为这样的 \(v \in \mathcal{L}(E)\) 全体所成的集合：\(x\) 是 \(v\) 的特征向量。则 A 是 \(\mathcal{L}(E)\) 的一个幺子代数。

代数 A 是全的：若 \(v \in A\) 在 \(\mathcal{L}(\mathrm{E})\) 中可逆，且 x 是 v 的对应于特征值 \(\lambda\) 的特征向量，则 \(\lambda \neq 0\) 且 \(v^{-1}(x) = \lambda^{-1}x\) ，这就证明了 \(v^{-1} \in A\) 。

代数 A 在 \(\mathcal{L}(\mathrm{E})\) 中是闭的。事实上，若 \((v_n)_{n\in \mathbf{N}}\) 是 A 中的序列，且 \(v_{n}\) 收敛于 \(v\in \mathcal{L}(\mathrm{E})\) ，则序列 \((\lambda_{n})_{n\in \mathbf{N}}\) （它满足 \(v_{n}(x) = \lambda_{n}x\) ）是有界的，故存在收敛于某个复数 \(\mu\) 的子序列，从而 \(v(x) = \mu x\) 。

假设 \(x\) 是 \(u\) 的特征向量且 \(u(x) = \lambda x\) 。代数 A 包含 \(u\) ，因而包含由 \(u\) 生成的闭幺全子代数 B，它是交换的。映射 \(\chi \colon \mathrm{B} \to \mathbf{C}\) 把 \(v\) 对应到 \(v\) 相对于 \(x\) 的特征值，它是 B 的一个特征标，且 \(\chi(u) = \lambda\) 。对每个 \(f \in \mathcal{O}(\mathrm{Sp}(u))\) ，由 I，第 75 页，命题 7，有 \(f(u) \in \mathrm{B}\) 且 \(\chi(f(u)) = f(\chi(u)) = f(\lambda)\) ，命题得证。

